\chapter{The Ledger}
\label{chap:the-ledger}

% EPISODE COVERAGE: Episodes 1 (completion) and 2 -- DISTINGUISHABLE through NUMERIC,
% Sample, ObservationProcess, BINARY, REPEATABLE, ComputationalProcess.closure
%
% AUDIENCE NOTE: Introduce the idea of a ledger as an accumulation of compiler
% agreements. Each typeclass is a new rule the compiler must accept. The ledger
% is irreversible: the compiler cannot un-accept a prior agreement.
%
% KEY CONCEPT TO ESTABLISH: ComputationalProcess.closure is the elaborator
% doing what measurement does -- closing an open account.
%
% TONE: Still accessible, but the Lean code starts to carry more weight.
% Use sidebars to protect non-technical readers.

\section{One Rule at a Time}
\label{se:one-rule-at-a-time}

% The typeclass as a compiler agreement.
% Introduce DISTINGUISHABLE:
%
%   class DISTINGUISHABLE (Value: Type) (Carrier: CarrierProcess Value) where
%     fact: Fact
%     symbol: Type Value
%     distinct? : Prop
%     different? : Type Value -> Prop
%     dec_distinct : DecidablePred different?
%
% Plain language: to be distinguishable is to be equipped with a method
% for telling things apart. The compiler will not call two things distinct
% until it has a decision procedure for distinguishing them.
%
% The type parameter pattern: every subsequent class is parameterized over
% (Value: Type) (Carrier: CarrierProcess Value). The compiler carries the
% context of the carrier forward into every new agreement.
%
% Sidebar: What is a typeclass? In Lean, a typeclass is a collection of
% operations that a type is required to provide. It is a contract. When
% you declare that something is DISTINGUISHABLE, you are promising the
% compiler that you can distinguish things of that type. The compiler
% holds you to the promise.

\section{The Chain of Admissibility}
\label{se:chain-of-admissibility}

% Walk through ADMISSIBLE, COUNTABLE, ENCODED, RESIDUE as a progressive
% tightening of the contract.
%
% ADMISSIBLE: adds counting_process and admissible? -- the first explicit
% admissibility predicate. A proposed continuation is admissible if it
% passes this test.
%
% COUNTABLE: adds index (IndexingProcess) and bounded? -- we can now
% place events on a number line and ask whether a range is finite.
%
% ENCODED: adds limit_process and encoding? -- we can compare sequences
% of events.
%
% RESIDUE: adds cauchy_process and representative? -- we can identify
% equivalence classes of sequences. The limit exists as a residue.
%
% Key metaphor: imagine a courtroom. ADMISSIBLE is the rule that says
% what evidence may be introduced. COUNTABLE is the rule that says how
% evidence must be numbered. ENCODED is the rule that says how two
% pieces of evidence may be compared. RESIDUE is the verdict: this
% evidence is equivalent to that evidence.
%
% Sidebar: Cauchy sequences. A Cauchy sequence is one where the terms
% get arbitrarily close to each other. It is the formal definition of
% "converging without needing to name the limit." The compiler accepts
% a Cauchy process as a proxy for a limit that might not exist in the
% current vocabulary.

\section{The First Multi-Fact Structure}
\label{se:first-multi-fact}

% Introduce Sample:
%
%   inductive Sample where
%     | initial_condition: Fact -> Limit -> Sample
%     | signal_response: Fact -> Limit -> Fact -> Limit -> Sample -> Sample
%
% A Sample is the first structure that holds more than one Fact. It
% models a stimulus and a response: an initial condition (one Fact, one
% Limit) and a signal_response that pairs a new Fact-Limit with a prior
% Sample.
%
% This is the ledger taking shape. Each signal_response appends a new
% witnessed event to the record. The record is ordered and irreversible
% by construction: you cannot build a signal_response without a prior Sample.
%
% The Limit type (from Episode 1) is an inductive sequence of Sequences.
% Together, Sample and Limit model the dual structure of observation:
% the thing observed (Limit, a sequence of encoded values) and the act
% of observation (Sample, a chain of Fact-witnessed events).

\section{The Clock and Its Complement}
\label{se:clock-and-complement}

% Episode 2: ObservationProcess and BINARY.
%
% ObservationProcess adds before: Limit and after: Limit -- the two
% sides of a witnessed transition.
%
% BINARY adds zero and one as Limits, and a bit as a Sample. The
% different? predicate has four cases: zero/zero, one/one (same),
% zero/one, one/zero (different). This is the compiler learning
% to count in binary.
%
% Connect to the covariant/contravariant duality from Chapter 1:
% the zero is the covariant ground state; the one is the contravariant
% step. The clock ticks T/F, T/F. Its complement ticks F/T, F/T.
% Together they tile every possible distinction.
%
% Sidebar: Why binary? Not because the universe is digital, but because
% a single distinction is the minimum unit of information. One bit. The
% compiler builds everything from yes/no decisions; of course the first
% number system it needs is binary.

\section{Closure}
\label{se:closure-first-appearance}

% The key arrival: ComputationalProcess.closure.
%
%   structure ComputationalProcess ... where
%     ...
%     closure: Study -> Study
%
% Show the full structure briefly. The closure field is a function
% from Study to Study -- it takes an open computational record and
% returns a closed one.
%
% What does "closed" mean? In the context of the compiler: a term
% is closed when all its free variables have been resolved, all its
% type holes have been filled, and every sub-expression has been
% type-checked. The elaborator has run to completion on the term.
%
% In the context of measurement: a ledger is closed when the proposed
% continuation is consistent with everything recorded. The measurement
% is complete.
%
% These are the same operation. ComputationalProcess.closure is the
% point where the book's central claim first becomes visible in the
% Lean code.
%
% NUMERIC follows: it wraps the closure in a class with a carrier
% (a Study), a lambda (the abstraction of the closure), and a
% related predicate (equality of closures).
%
% End of chapter: the compiler has now agreed to twelve rules. It can
% distinguish, count, encode, take limits, observe, work in binary,
% repeat experiments, and close. The ledger is not yet full. But the
% compiler knows what a closure is.

